\section{Постановка}
\label{sec:problem}
\porozhdeno{\Claim{P01}{Насколько меняются таксономические назначения и профили одних почвенных образцов при выборе готовой модели классификатора? Мы проводим описательное сравнение моделей, связанных с SILVA, GTDB и Greengenes2 (далее GG2), при одном алгоритме \texttt{classify-sklearn} и фиксированных входных последовательностях. Измеряются покрытие назначениями, буквальное совпадение родословных и расстояние профилей; задача не предусматривает выбора наиболее точной модели без эталона.}}

\porozhdeno{\Claim{P02}{Набор содержит \ProjectSamplesSoil{} почвенных образцов Крылатского и Люблина и отрицательный контроль K\_S63 --- амплифицированную воду. Контроль обработан вместе с почвами, но исключён из почвенных сводок. В таблице после денойзинга \ProjectAsvAll{} вариантов ампликонной последовательности (ASV), из них \ProjectAsvSoil{} встречаются в почвах. Единица профильного сравнения --- образец, а единица сравнения назначений --- ASV; разные ASV и технические исполнения не считаются независимыми биологическими повторностями.}}

\porozhdeno{\Claim{P03}{Правило сравнения зафиксировано в методике 10 ноября 2026 года до сравнительного расчёта; постановка и выбор представленных результатов оформлены 18 ноября. Поздние наблюдения ниже относятся к описательному разбору данных, а не к заранее зарегистрированным гипотезам.}}
